% from ega3-7-fr.tex
\oldpage[III]{208}
avec $Y'=\operatorname{Spec}(f_*(\mathscr O_X))$ ; le morphisme fini $g$ est alors
tel que $g_*(\mathscr O_{Y'})=f_*(\mathscr O_X)$ soit localement libre au
voisinage de $y$, et sa fibre en $y$ est le spectre d'une algèbre séparable sur
$\kappa(y)$ (II, 1.5.1). Nous en déduirons au chap. IV qu'il y a un voisinage
ouvert $U$ de $y$ dans $Y$ tel que pour la restriction $g^{-1}(U)\to U$ de $g$,
\emph{toute} fibre $g^{-1}(z)$ (où $z\in U$) soit spectre d'une algèbre séparable
sur $\kappa(z)$ (c'est ce que nous appellerons un \emph{revêtement étale} de
$U$) ; il résultera alors de (7.8.7, (ii)) que l'hypothèse faite sur le point
$y$ dans (7.8.6) est vérifiée aussi en tous les points d'un voisinage de $y$.

\medskip
\textnormal{(ii)} Nous verrons au chap. V que, même si $X$ est projectif sur $Y$
(et même s'il est en outre « simple » sur $Y$, propriété qui sera définie au
chap. IV), le $\mathscr O_U$-Module $\mathscr Q$ de (7.8.9) n'est pas
nécessairement localement libre ; en d'autres termes, $\mathscr O_X$ (sous ces
conditions) n'est pas nécessairement cohomologiquement plat en \emph{dimension
1} sur $Y$ au point $y$. Au chap. V, nous interpréterons $\mathscr Q$ comme le
faisceau des 1-différentielles du schéma de Picard de $X$ par rapport à $Y$ le
long de la section unité.
\end{remarks}

\subsection{Application aux morphismes propres : III. Invariance de la
caractéristique d'Euler-Poincaré et du polynôme de Hilbert}
\label{subsection:III.7.9.fr}
\label{III.7.9.fr}

\begin{env}[7.9.1]
\label{III.7.9.1.fr}
Soient $A$ un anneau, $M$ un $A$-module projectif de type fini ; rappelons
(Bourbaki, \emph{Alg. comm.}, chap. II, \S~5, n\textsuperscript{o}~2) qu'il
revient au même de dire que le $\mathscr O_X$-Module associé $\widetilde M$ sur
$X=\operatorname{Spec}(A)$ est \emph{localement libre de type fini}. Pour tout
$\mathfrak p\in\operatorname{Spec}(A)$ on appelle \emph{rang de $M$ en
$\mathfrak p$} et on note $\operatorname{rang}_{\mathfrak p}(M)$ le rang du
$A_{\mathfrak p}$-module libre $M_{\mathfrak p}$ (ou encore le rang en
$\mathfrak p$ du $\mathscr O_X$-Module localement libre $\widetilde M$). On a
donc
\[
  \operatorname{rang}_{\mathfrak p}M
  =\operatorname{rang}_{\mathfrak p}(M_{\mathfrak p})
  =\operatorname{rang}_{\kappa(\mathfrak p)}
  \bigl(M\otimes_A\kappa(\mathfrak p)\bigr).
  \tag{7.9.1.1}\label{III.7.9.1.1.fr}
\]
\end{env}

\begin{proposition}[7.9.2]
\label{III.7.9.2.fr}
Soit $P_\bullet$ un complexe fini de $A$-modules projectifs de type fini, et
pour tout $A$-module $M$, soit
$T_\bullet(M)=H_\bullet(P_\bullet\otimes_A M)$. Alors, pour tout
$\mathfrak p\in\operatorname{Spec}(A)$, on a
\[
  \sum_i(-1)^i\operatorname{rang}_{\kappa(\mathfrak p)}
  T_i\bigl(\kappa(\mathfrak p)\bigr)
  =\sum_i(-1)^i\operatorname{rang}_{\mathfrak p}(P_i).
  \tag{7.9.2.1}\label{III.7.9.2.1.fr}
\]
\end{proposition}

\begin{proof}
En effet, on a par définition
$T_i(\kappa(\mathfrak p))=H_i(P_\bullet\otimes_A\kappa(\mathfrak p))$ et,
compte tenu de (7.9.1.1), la formule (7.9.1.2) n'est autre que l'invariance de
la caractéristique d'Euler-Poincaré d'un complexe fini d'espaces vectoriels de
dimension finie par passage à l'homologie ($0_{\mathrm I}$, 11.10.2).
\end{proof}

\begin{corollary}[7.9.3]
\label{III.7.9.3.fr}
La fonction
\[
  \mathfrak p\longmapsto
  \sum_i(-1)^i\operatorname{rang}_{\kappa(\mathfrak p)}
  T_i\bigl(\kappa(\mathfrak p)\bigr)
\]
est localement constante dans $\operatorname{Spec}(A)$.
\end{corollary}

\begin{theorem}[7.9.4]
\label{III.7.9.4.fr}
Soient $Y$ un préschéma localement noethérien, $f:X\to Y$ un morphisme propre,
$\mathscr P_\bullet$ un complexe fini de $\mathscr O_X$-Modules cohérents et
$Y$-plats. Si l'on pose
\[
  \mathscr T_\bullet(\mathscr M)
  =\mathscr H^{-\bullet}
  \bigl(f,\mathscr P^\bullet\otimes_{\mathscr O_Y}\mathscr M\bigr)
  \qquad\text{(cf. (7.7.1.1))}
\]
la fonction

